% Imported from LaTeX_and_pdfs/AGAT_dump.tex; see provenance.json.
\part{Category Theory and Databases}
 
	
\section{Types}

\subsection{Data models for nested arrays, dictionaries, and tabular data}

\begin{definition}[Container types]
	\textbf{Container types} - arrays, key/value pair dictionaries (or: hashes, association lists)
\end{definition}

\begin{definition}[atomic types]
	basic atomic types (e.g. numbers, strings, Booleans)
\end{definition}

\begin{definition}[nesting]
	\textbf{nesting}: containers may contain atomic values as well as other containers.
\end{definition}

\begin{definition}[flat]
	Tabular data model (i.e. tables) is \textbf{flat}: field contains \textbf{atomic} values.
\end{definition}
cf. Number 2, Data Models and Languages, Grust and Duta (2017) \cite{agat:GrDu2017}


\subsection{Typed Data, Untyped data, in the relational data model}

\begin{definition}
	Untyped data models - text, JSON, and tabular data models (e.g. CSV) do \textbf{not} enforce values (container or atomic) to be of specific types. 
	
	These data models are thus referred to as being \textbf{untyped}.
\end{definition}
cf. Grust and Duta (2017) \cite{agat:GrDu2017}


cf. Number 4, "The Relational Data Model", Grust and Duta (2017) \cite{agat:GrDu2017}

\begin{definition}[Types]
Let $\mathbb{T} \equiv $ set of all data types (built-in and user-defined).  \\
$\forall \, $ value $v\in  \mathbb{V}$ stored in a relation cell must be of type $t\in \mathbb{T}$. \\
e.g. When PostgreSQL starts, $\mathbb{T}$ initialized as 
\begin{equation}
 \mathbb{T} = \lbrace \texttt{boolean}, \texttt{integer}, \texttt{text}, \texttt{bytea}, \dots \rbrace 
\end{equation}
Consider category \textbf{Text} s.t. $\text{Obj}(\mathbf{Types}) \equiv DT \in \text{Obj}(\textbf{Set})$. Then $\mathbb{T} \equiv DT$, denoting \emph{data type}.
\end{definition}

\begin{definition}[Values]
$\forall \, v \in \mathbb{V}$ stored in a relation cell, $v$ is an element of the set of all values $\mathbb{V}$. 

in the relational data model, all values $v\in \mathbb{V}$ are "atomic."

\[
\mathbb{V} = \lbrace \texttt{true}, \texttt{false}, 0, -1, 1, -2, 2 , \dots \rbrace
\]	

Here, I'll use the notation $V$ to denote $\mathbb{V}$, the set of all values.
\end{definition}

\begin{definition}[Domains]
	$\forall \, t \in DT$, its domain $\text{dom}(t) := $ set of all values of type $t$ (i.e. $\text{dom}(\cdot)$ is a function with signature 
	\[
	DT \to 2^V
	\])
	e.g. $\text{dom}(\texttt{integer}) = \lbrace 0 , -1, 1, -2, 2, \dots \rbrace $ \\
	$\text{dom}(\texttt{boolean}) = \lbrace \texttt{true}, \texttt{false} \rbrace$ 
\end{definition} 

\begin{definition}[type specification]
	type specification $:= $ function $\pi : U \to \mathbf{DT}$ (Spivak's notation) $ \equiv \pi : U \to DT$, $U, DT \in \textbf{Set}$.
	
	set $DT \equiv $ set of data types for $\pi$, set $U \equiv $ domain bundle for $\pi$.
	
	$\forall \, t \in DT$, preimage $\pi^{-1}(t) \subset U$, \, $\pi^{-1}(t) \equiv $ domain of $t$, $x \in \pi^{-1}(t) \equiv $ object of type $T$.

	To reconcile Grust and Duta (2017) \cite{agat:GrDu2017}'s definition of types above, use this notation:
	\begin{equation}
	\boxed{
	\begin{gathered}
	\textbf{ type specification } \begin{aligned} & \quad \\ 
	& \pi : V \to DT \\ 
	& \pi(v) \in DT \end{aligned} \\
	\text{dom} \equiv \pi^{-1}: DT \to V, \text{ i.e. } \\
	\pi^{-1}(t) \subset t \qquad \, \forall \, t \in DT  
	\end{gathered}}
	\end{equation}
\end{definition}
	cf. Spivak (2009) \cite{agat:Spiv2009}

\begin{corollary}[type specification]
	If $v$ has type $T$, $\pi(v) = T$, $\Longrightarrow v \in \pi^{-1}(T)$
\end{corollary}

\begin{proposition}[\texttt{CREATE DOMAIN}]
	Consider new type $t'$ (SQL command \texttt{CREATE DOMAIN}) so 
	\begin{equation}
	\begin{gathered}
	t' \in DT \\
	\pi^{-1}(t') \subseteq \pi^{-1}(t)
	\end{gathered}
	\end{equation}
\end{proposition}


\section{Relational Data Model}

	\textbf{relational data model} maybe understood as a typed variant of the tabular data model. 
	\begin{enumerate}
		\item $\exists \, $ only 1 container type: \textbf{table} (or: multisets) \textbf{ of rows}
		\item all rows are of same \textbf{row type} which is declared when table is created.
		\item row type \textbf{consists} of sequence of \textbf{atomic types}.
	\end{enumerate}

In the relational data mode, data is exclusively organized in \textbf{relations}, i.e. sets of tuples of data. Data in each \textbf{attribute} (tuple component) is \textbf{atomic}, and of declared \textbf{type}.

\subsection{Schemata and Relations}

In the relational data model, each \textbf{attribute} of a table has a \textbf{declared type}.

If attribute has declared type $t$, the RDBMS will exclusively store values $v$ in that attribute s.t. 
\begin{enumerate}
	\item $v \in \text{dom}(t)$ i.e. $v \in \pi^{-1}(t)$
	\item $v$ can successfully be casted to type $t$
\end{enumerate}

\subsection{Attributes (Columns)}

Let $\mathbb{A}$ denote set of attribute names of all relations.

\subsection{Attribute types}
$\forall \, $ attribute $a\in \mathbb{A}$ has declared (attribute) type $\text{type}(a) = t\in DT$ (i.e. $\text{type}(\cdot)$ is a function with signature $\mathbb{A} \to \mathbb{T}$).

Consider Definition 2.2.3 of Spivak (2009) \cite{agat:Spiv2009}, 

\begin{definition}[simple schema of type $\pi$, $(C,\sigma)$]
	Let type specification $\pi :V \to DT$. 
	
	\textbf{simple schema} of type $\pi$ consists of pair $(C, \sigma)$ where $C$ is a finite (totally) ordered set and function $\sigma : C \to DT$.  $C \equiv $ \textbf{column set} or \textbf{set of attributes} for $\sigma$ and $\pi$ as \emph{type specification for} $\sigma$.
\end{definition}

Compare the notation above. Conclude that
\[
\begin{gathered}
C \equiv \mathbb{A}, \quad \, c \equiv a \\
\sigma(c) \equiv \text{type}(a)
\end{gathered}
\]

Here, choose the following notation:
\[
\boxed{\begin{gathered}
\sigma : C \to DT \\
\pi^{-1}: DT \to V \\
\pi^{-1} \circ \sigma \equiv \text{val} \equiv \text{set of (admissable) attribute values for attribute (column) $a$ ($c$)}
\end{gathered} }
\]
Also, summarize our definitions with these commutative diagrams:
\begin{equation}
\begin{gathered}
\begin{tikzpicture}
\matrix (m) [matrix of math nodes, row sep=3em, column sep=4em, minimum width=2em]
{
	V     \\
	DT  \\
};
\path[->]
(m-1-1) edge node [left] {$\pi$} (m-2-1)
;
\end{tikzpicture}  \qquad \, 
\begin{tikzpicture}
\matrix (m) [matrix of math nodes, row sep=3em, column sep=4em, minimum width=2em]
{
	v     \\
	t  \\
};
\path[|->]
(m-1-1) edge node [left] {$\pi$} (m-2-1)
;
\end{tikzpicture}  \qquad \, 
\begin{tikzpicture}
\matrix (m) [matrix of math nodes, row sep=3em, column sep=4em, minimum width=2em]
{
	V &  \pi^{-1}(DT)  \\
	DT &  \\
};
\path[left hook->]
(m-1-2) edge node [above] {$i$} (m-1-1)
;
\path[->] 
(m-2-1) edge node [left] {$i \circ \pi^{-1}$} (m-1-1)
edge node [left] {$\pi^{-1}$} (m-1-2)
;
\end{tikzpicture} \qquad \, 
\begin{tikzpicture}
\matrix (m) [matrix of math nodes, row sep=3em, column sep=4em, minimum width=2em]
{
	V &  V_t \subseteq V  \\
	t &  \\
};
\path[left hook->]
(m-1-2) edge node [above] {$i$} (m-1-1)
;
\path[|->] 
(m-2-1) edge node [left] {$\pi^{-1}$} (m-1-2)
;
\end{tikzpicture} \\
\begin{tikzpicture}
\matrix (m) [matrix of math nodes, row sep=3em, column sep=4em, minimum width=2em]
{
	V &  \pi^{-1}(DT)  \\
	DT &  C \\
};
\path[left hook->]
(m-1-2) edge node [above] {$i$} (m-1-1)
;
\path[->] 
(m-2-1) edge node [left] {$i \circ \pi^{-1}$} (m-1-1)
edge node [left] {$\pi^{-1}$} (m-1-2)
(m-2-2) edge node [right] {$\pi^{-1} \circ \sigma \equiv \text{dom} \circ \text{type} $} (m-1-2)
edge node [above] {$\sigma$} (m-2-1)
;
\end{tikzpicture} \qquad \, 
\begin{tikzpicture}
\matrix (m) [matrix of math nodes, row sep=3em, column sep=4em, minimum width=2em]
{
	V &  \pi^{-1}(\sigma(c_i)) = V_{c_i}  \\
	\sigma(c_i) \equiv t_i & c_i \\
};
\path[left hook->]
(m-1-2) edge node [above] {$i$} (m-1-1)
;
\path[|->] 
(m-2-1) edge node [left] {$\pi^{-1}$} (m-1-2)
(m-2-2) edge node [above] {$\sigma$} (m-2-1)
edge node [right] {$\pi^{-1} \circ \sigma$} (m-1-2)
;
\end{tikzpicture} 
\end{gathered}
\end{equation}


\begin{definition}[Relation Schema]
	A relation schema associates a relation (table?) name $R$ with its set of declared attributes (a subset of $\mathbb{A}$)
	\begin{equation}
	(R, \lbrace a_1, \dots a_n \rbrace)
	\end{equation}
Common notation: $R(a_1, \dots a_n)$, so that $R$ is called a $n$-ary relation.

More notation: $\text{sch}(R) = \lbrace a_1, \dots a_n \rbrace$, and $\text{deg}(R) = n$ (degree).
		
\textbf{Relational database schema}: a non-empty finite set of relation schemata makes a \textbf{relational database schema}
\begin{equation}
\lbrace (R_1, \alpha_1), (R_2, \alpha_2) , \dots \rbrace 
\end{equation}
where $\alpha_i \subseteq \mathbb{A}_i$. In a relational database schema, the relation names $R_i$ are unique.		
\end{definition}

\begin{definition}[Tuple]
	Given relation (i.e. table) $R(a_1, \dots , a_n)$, a \textbf{tuple} $t$ of $R$ maps attributes to values, i.e. $t$ is a function with signature $\lbrace a_1, \dots a_n \rbrace \to V$ with 
	\[
	\forall \, a \in \lbrace a_1 , \dots , a_n \rbrace : t(a) \in \text{val}(a)
	\]
	Common notation for $t(a)$ is $t.a$.
\end{definition}

Recall that $\text{val}(a) := \text{dom}(\text{type}(a))$.

Take note that \textbf{tuple}, defined by Grust and Duta (2017) \cite{agat:GrDu2017}, is the \emph{same} as \textbf{record} or \textbf{row} $r$, defined by Spivak (2009) \cite{agat:Spiv2009}.

\begin{definition}[record or row]
	\textbf{record} or \textbf{row} on $(C, \sigma)$
	\begin{equation}
	\begin{aligned}
	& r: C \to V_{\sigma} \equiv V_c \\ 
	& r(c) \equiv v_{rc} \in \pi^{-1}\circ \sigma(c)
	\end{aligned}
	\end{equation}
\end{definition} 

SQL \texttt{CREATE TABLE} command prescribes an order of the attributes of a relation. This deviates from relational data model's tuple model (name-to-value mapping).

\begin{definition}[Row]
	Given SQL table $R(a_1, \dots, a_n) \equiv R(\alpha) \equiv \tau$, a row $r$ of $\tau$ is an ordered sequence ($a_i$ is called the $i$th column)
	\begin{equation}
	r = (v_1, \dots, v_n) \in \text{val}(a_1) \times \dots \times \text{val}(a_n) \equiv \pi^{-1} \circ \sigma(a_1) \times \dots \times \pi^{-1} \circ \sigma(a_n)
	\end{equation}
	Thus, $r$ is a function $\lbrace 1, \dots , n \rbrace \to V$ with $\forall \, i \in \lbrace 1 , \dots , n \rbrace$, $r(i) \in \pi^{-1} \circ \sigma(a_i)$
\end{definition}

The set of tuples (rows) stored in a relation (table) is expected to change frequently.

\begin{definition}[Relation instance (state)]
	The current finite set of tuples $t_i \equiv r_i$ of relation (table) $R(a_1, \dots a_n) \equiv (A, \sigma)$ is called the relation's \textbf{instance} (or \textbf{state}).

\[
\text{inst}(R) = \lbrace t_1 , t_2 , \dots t_m \rbrace \Longrightarrow \equiv \Gamma^{\pi}(\sigma) \equiv \Gamma(A, \sigma)
\]

Database (instance) state - The database instance comprises instances of all its relations.
\[
= \lbrace \Gamma(A^i, \sigma^i) \rbrace_i
\]
\end{definition} 

\subsection{Constraints}

cf. Number 5, "Constraints" of Grust and Duta (2017) \cite{agat:GrDu2017}.

\begin{definition}[Constraints]
	An \textbf{integrity constraints} specifies conditions which table states have to satisfy at all times. 
	
	Current set of constraints, $\mathbb{C}$, is integral part of database schema:
	\[
	(\lbrace (R_1, \alpha_1), (R_2, \alpha_2), \dots  \rbrace, \mathbb{C})
	\]
	Set of constraints $\mathbb{C} \equiv$ set of morphisms of table $\tau \equiv \lbrace (R_1,\alpha_1), (R_2, \alpha_2), \dots \rbrace $, $C \subset \text{Mor}{\tau}$. 
\end{definition} 

RDBMS will refuse table state changes that violate any constraint $c\in C$.

\subsection{Key constraints}

\begin{definition}[Key]
Key of a table $R(a_1, \dots, a_n) \equiv (\alpha, \sigma) \equiv (\alpha, \text{type}) :=$ set of columns $K \subseteq \lbrace a_1, \dots, a_n \rbrace$ that unique identifies rows of $R$:
\[
\forall \, t, u \in \text{inst}(R), \, t.K = u.K \Longrightarrow t = u
\]
Read: "If 2 rows agree on the columns in $K$, they are indeed the same row."

I will give the following change of notation a try:
\begin{equation}
\forall \, r_i , r_j \in \lbrace r_i \rbrace_i^{\tau}, \, \text{ if } r_i(K) = r_j(K), \text{ then } r_i = r_j
\end{equation}
\end{definition} 

\begin{equation}
\begin{gathered}
\begin{tikzpicture}
\matrix (m) [matrix of math nodes, row sep=3em, column sep=4em, minimum width=2em]
{
	A &  DT  \\
	V_A &  \\
};
\path[->] 
(m-1-1) edge node [above] {$\sigma$} (m-1-2)
edge node [left] {$r$} (m-2-1)
(m-2-1) edge node [right] {$\pi^{-1}$} (m-1-2)
;
\end{tikzpicture} \qquad \, 
\begin{tikzpicture}
\matrix (m) [matrix of math nodes, row sep=3em, column sep=4em, minimum width=2em]
{
	a &  t  \\
	V_A &  \\
};
\path[|->] 
(m-1-1) edge node [above] {$\sigma$} (m-1-2)
edge node [left] {$r$} (m-2-1)
(m-2-1) edge node [right] {$\pi^{-1}$} (m-1-2)
;
\end{tikzpicture}
\end{gathered}
\end{equation}


Here's a "dictionary" between the definitions so far for relational databases and familiar terms for tables filled with data: cf. 08-30, Week 1 slides of Yang (2012) \cite{agat:Yang2012}.

database - collection of relations (or tables) \\
relation - table \\
attributes - columns \\
tuple - row - record \\
relation schema - heading (heading for a table?, Grust and Duta (2017) \cite{agat:GrDu2017}) \\
relation contents - body (body of a table?, Grust and Duta (2017) \cite{agat:GrDu2017})

\begin{tabular}{l c r}
	CSV & Relational Model & SQL \\
	& Domain & Domain \\ 
	& Type & Type \\
	& Schema & Schema \\
	File & Relation & Table \\
	Line & Tuple & Row \\
	Field & Attribute & Column
\end{tabular}

Indeed, even Grust and Duta (2017) \cite{agat:GrDu2017} remarks: "You will find that textbooks, papers, practitioners, academics, these slides, and even PostgreSQL use a mixture of terminology. Deal with it."


cf. Slides Number 7, "Referential Integrity" of Grust and Duta (2017) \cite{agat:GrDu2017}.

\begin{definition}[Foreign Keys]
	Let $(S, \alpha)$ and $(T, \beta)$ denote 2 relational schemata (not necessarily distinct), \\
	where $K = \lbrace b_{j_1}, \dots , b_{j_k} \rbrace \subseteq \beta$ is the primary key of $T$. \\
	Let $F = \lbrace a_{i_1}, \dots , a_{i_k} \rbrace \subseteq \alpha $ with $\text{type}(a_{i_h}) = \text{type}(b_{j_h})$, $h \in 1, \dots, k$. 
	
	$F$ is a \textbf{foreign key} in $S$ \textbf{referencing} $T$, if 
	\[
	\forall \, s \in \text{inst}(S) : \exists \, t \in \text{inst}(T) : s.F = t.K
	\]
\end{definition} 
The $\forall$ and $\exists $ condition validates the assumption of $K$ being a key in target $T$, i.e. that there \textbf{exists a row} in table $T$ whose $K$ identifier matches that of table $S$

$K$ being a key in target $T$ validates the assumption that there is \textbf{no more than one row} of in a table $T$ with a matching key $K$.

In general, \textbf{ a foreign key $F$ is not a key} in source table $S$. 2 rows $s_1, s_2 \in \text{inst}(S)$ with $s_1.F = s_2.F$ can refer to the same row in target $T$.

Here is a great example of a foreign key: \url{https://www.w3schools.com/sql/sql_foreignkey.asp}

So from w3schools, 

"A FOREIGN KEY is a key used to link 2 tables together.

A FOREIGN KEY is a field (or collection of fields) in 1 table that refers to the PRIMARY KEY in another table.

The table containing the foreign key is called the child table, and the table containing the candidate key is called the referenced or parent table.

"Persons" table:

\begin{tabular}{|l |c |c |r|}
\textbf{PersonID} & \textbf{LastName} & \textbf{FirstName} & \textbf{Age} \\
1 & Hansen & Ola & 30 \\
2 & Svendson & Tove & 23 \\
3 & Pettersen & Kari & 20 
\end{tabular}

"Orders" table:

\begin{tabular}{ | l | c | r | }
	\textbf{OrderID} & \textbf{OrderNumber} & \textbf{PersonID} \\ 
	1 & 77895 & 3 \\
	2 & 44678 & 3 \\
	3 & 22456 & 2 \\
	4 & 24562 & 1 
\end{tabular}

Notice that the "PersonID" column in the "Orders" table points to the "PersonID" column in the "Persons" table.

The "PersonID" column in the "Persons" table is the PRIMARY KEY in the "Persons" table.

The "PersonId" column in the "Orders" table is a FOREIGN KEY in the "Orders" table.

\quad \\
Thus, in other words, \\
given 2 tables (i.e. 2 relational schemata, not necessarily distinct), $(S, \alpha)$, $(T,\beta)$, \\
given primary key $PK = \lbrace b_{j_1}, \dots, b_{j_k} \rbrace \subseteq \beta$, primary key of of $T$, and \\
given foreign key $FK = \lbrace a_{i_1} , \dots , a_{i_k} \rbrace \subseteq \alpha$ s.t.
\[
\begin{gathered}
\text{type}(a_{i_h}) = \text{type}(b_{j_h}), \quad \, h \in 1, 2 , \dots k \\
\text{(alternate notation)}  \sigma(a_{i_h}) = \sigma(b_{j_h})
\end{gathered}
\]
then \\
$\forall \, s \in \text{inst}(S), \, \exists \, t \in \text{inst}(T)$ ((alternate notation) $s\in \Gamma(S), \, t\in \Gamma(T)$), s.t.
\[
s(FK) = t(PK)
\]

Describe this categorically as
\[
\begin{tikzpicture}
\matrix(m)[matrix of math nodes, row sep=4em, column sep=4em]
{
	S   &  T  \\
};
\path[->]
(m-1-1) edge node[auto]{$FK$} (m-1-2)
;
\end{tikzpicture} 
\]
which is from the introductory talk by Spivak (2012)\cite{agat:Spiv2012}, "What does equivalence of paths mean?", where it states, "arrows represent foreign keys", "from table $a$ to table $b$," and where "we can take any record in table $a$ and return a record in table $b$.

\subsection{Functional Dependency}

\begin{definition}[Functional Dependency (FD)]
	cf. Slides Number 10, "Functional Dependencies" of Grust and Duta (2017) \cite{agat:GrDu2017}.
	
	Let $(R, \alpha) \equiv $ relational schema. Given $\beta \subseteq \alpha$, $c\in \alpha$, \\
	\textbf{functional dependency} $\beta \to c$ holds in $R$ if 
	\[
	\forall \, t, u \in \text{inst}(R), \, t.\beta = u.\beta \Longrightarrow t.c = u.c
	\]
	Notation: the functional dependency $\beta \to \lbrace c_1, \dots , c_n \rbrace $ abbreviates set of FDs $\beta \to c_1, \dots , \beta \to c_n$.
	
	cf. \href{https://en.wikipedia.org/wiki/Functional_dependency}{"Functional dependency", wikipedia}
	
	A set of attributes $X \subseteq R$ ($\equiv X \subseteq \alpha$ for $R = (R,\alpha)$) is said to functionally determine another set of attributes $Y \subseteq R$ ($\equiv Y \subseteq \alpha$), written $X\to Y$. ($\equiv \beta \to c$) \\
	iff $\forall \, X$ value in $R$ is associated with precisely 1 $Y$ value in $R$; $R$ is then said to satisfy the functional dependency $X\to Y$
\end{definition}

e.g. Cars: $\beta = $ vehicle identification number (VIN), $c=$ Engine capacity (because assume a car cannot have 2 engines).

e.g Employee department: \\
$\beta = $ employee ID, \qquad \, $c = $ employee name. \\
$\beta = $ employee ID, \qquad \, $c = $ department ID. \\
$\beta = $ department ID, \qquad \, $c = $ department name.

Recall definition of injective function, $f:X \to Y$,
\[
\forall \, a,b \in X, \, f(a) = f(b) \Longrightarrow a =b
\]

Instead of writing the functional dependency (FD) with this notation: given $\beta, c \subseteq \alpha$, 
\[
\begin{gathered} 
\forall \, r_1, r_2 \in \text{inst}(R) \equiv \Gamma(\alpha), \, r_1(\beta) = r_2(\beta), \Longrightarrow r_1(c) = r_2(c) \\
\end{gathered} 
\]
\emph{rewrite} this as follows:

\begin{equation}
\begin{gathered} 
	\text{ Given } \\
	\begin{aligned} 
	& \beta : \Gamma(\alpha) \to \text{val}(\beta) \equiv \text{dom}(\text{type}(\beta)) \equiv \pi^{-1} \circ \sigma(\beta), \, \\
	& c : \Gamma(\alpha) \to \text{val}(c) \equiv \text{dom}(\text{type}(c)) \equiv \pi^{-1} \circ \sigma(c), \end{aligned} \\
\forall \, r_1, r_2 \in \text{inst}(R) \equiv \Gamma(\alpha), \, \beta(r_1), \beta(r_2) \in \text{val}(\beta)  \\
\text{ Then } \beta \in \text{Hom}(\text{inst}(R), \text{val}(\beta)) \equiv \text{Hom}(\Gamma(\alpha), \pi^{-1}\circ \sigma(\alpha))
\end{gathered} 
\end{equation}

So define a \textbf{functional dependency} $FD$ as 
\begin{equation}
\begin{gathered}
	\begin{aligned} 
&	\text{FD} : 2^{\alpha} \to 2^{\alpha} \\
& \text{FD} : \beta \mapsto c \text{ i.e. } \boxed{ FD(\beta) = c }
\end{aligned} 
\end{gathered}
\end{equation}
where $\beta, c \in 2^{\alpha} \equiv $ power set of $\alpha$, i.e. set of all subsets of $\alpha$, so that $\beta, c \subseteq \alpha$.  

Let $FD^{-1}:2^{\alpha} \to 2^{\alpha}$, s.t. $FD^{-1}(c) = \beta$.

So if $\forall \, r_1, r_2 \in \text{inst}(R) \equiv \Gamma(\alpha)$, for $\beta \subseteq \alpha$, $\beta(r_1), \beta(r_2) \in \text{val}(\beta)$. Then $\beta \in \text{Hom}(\text{inst}(R), \text{val}(\beta)) \equiv \text{Hom}(\Gamma(\alpha), \pi^{-1}\circ \sigma(\alpha))$.

Suppose $\beta(r_1) = \beta(r_2)$. The \emph{key insight} is the following:

\begin{equation}
\beta(r_1) = \beta(r_2) = \text{FD}^{-1}(c)(r_1) = \text{FD}^{-1}(c)(r_2) = (\text{FD}^{-1}\circ c)(r_1) = (\text{FD}^{-1}\circ c)(r_2)
\end{equation}

Then $c(r_1) = c(r_2)$ implies that $\text{FD}^{-1}$ is \emph{injective}.

Then $\text{FD}^{-1}$ is a \emph{monomorphism}, and $\text{FD}^{-1}$ is injective, since $\forall \, r_1, r_2 \in \text{inst}(R) \equiv \Gamma(\alpha)$,
\begin{equation}
\beta(r_1) = \beta(r_2) =  (\text{FD}^{-1}\circ c)(r_1) = (\text{FD}^{-1}\circ c)(r_2) \Longrightarrow c(r_1) = c(r_2)
\end{equation} 

\[
\begin{tikzpicture}
\matrix (m) [matrix of math nodes, row sep=1.8em, column sep=1.8em, minimum width=1.2em]
{
	\text{inst}(R) & \text{val}(C) & \text{val}(\beta) \\
};
\path[->]
(m-1-1) edge node [above] {$c$} (m-1-2)
(m-1-2) edge node [above] {$\text{FD}^{-1}$} (m-1-3)
;
\end{tikzpicture} 
\]

Then conclude that 
\begin{theorem}
	Given a relational schema $(R, \alpha)$, a set of attributes $\alpha$, instance of the relation (i.e. a table), $\text{inst}(R) \equiv \Gamma(\alpha)$, and $\beta, c \subseteq \alpha$, \\ for morphisms $\beta, c$, 
\[	\begin{aligned}
		& \beta \in \text{Hom}(\Gamma(\alpha), \text{val}(\beta)) \\
		& c \in \text{Hom}(\Gamma(\alpha), \text{val}(c))
	\end{aligned}
	\]
Then for $\text{FD}:2^{\alpha} \to 2^{\alpha}$, $\text{FD}(\beta) =c$, $\text{FD}^{-1}$ is a \textbf{monomorphism}, i.e. $\forall \, r_1, r_2 \in \Gamma(\alpha)$,
\[
\beta(r_1) = \beta(r_2) = (\text{FD}^{-1}\circ c)(r_1) = (\text{FD}^{-1}\circ c)(r_2) \text{ i.e. } (\text{FD}^{-1}\circ c) = (\text{FD}^{-1} \circ c)
\]
then $c(r_1) = c'(r_2) $ i.e. $c=c'$.
\end{theorem}

Database Queries and Constraints via Lifting Problems. David I. Spivak.
\url{https://arxiv.org/pdf/1202.2591.pdf}

\section{Databases and Categories}

cf. Spivak (2012) \cite{agat:Spiv2012}.

$A \xrightarrow{\text{FK}} B$, \, $A, B \in \text{Obj}(\mathbf{DB})$


